% special-elements.tex

%%%%%%%%%%%%%%%
\begin{frame}{}
  \begin{definition}[极大元与极小元 (Maximal/Minimal)]
    令 $\preceq \;\subseteq X \times X$ 是 $X$ 上的偏序。
    设 $a \in X$,\\[5pt]

    \pause
    如果
    \[
      \forall x \in X\; \lnot \big(a \prec x \big),
    \]
    则称 $a$ 是 $X$ 的\red{极大元}。\\[10pt]

    \pause
    如果
    \[
      \forall x \in X\; \lnot \big(x \prec a \big),
    \]
    则称 $a$ 是 $X$ 的\red{极小元}。
  \end{definition}

  \pause
  \vspace{0.30cm}
  \begin{center}
    \red{$Q:$ 极大/极小元是否一定存在? 如果存在, 是否唯一?}
  \end{center}
\end{frame}
%%%%%%%%%%%%%%%

%%%%%%%%%%%%%%%
\begin{frame}{}
  \begin{exampleblock}{}
    \[
      (\R, \le)
    \]
    \pause
    \begin{center}
      无极大元、无极小元
    \end{center}
  \end{exampleblock}

  \pause
  \vspace{0.30cm}
  \begin{exampleblock}{}
    \[
      (\N, \le)
    \]
    \pause
    \begin{center}
      无极大元、有唯一极小元 $0$
    \end{center}
  \end{exampleblock}
\end{frame}
%%%%%%%%%%%%%%%

%%%%%%%%%%%%%%%
\begin{frame}{}
  \begin{center}
    \fig{width = 0.40\textwidth}{figs/Hasse-Powerset}
  \end{center}
\end{frame}
%%%%%%%%%%%%%%%

%%%%%%%%%%%%%%%
\begin{frame}{}
  \fig{width = 0.60\textwidth}{figs/powerset}
\end{frame}
%%%%%%%%%%%%%%%

%%%%%%%%%%%%%%%
\begin{frame}{}
  \begin{exampleblock}{}
    \[
      (\set{2, 3, 4, \dots} = \N \setminus \set{0, 1}, \mid)
    \]
  \end{exampleblock}

  \pause
  \fig{width = 0.40\textwidth}{figs/infinite-divisors}

  \pause
  \begin{center}
    无极大元、有无穷多个极小元 (所有的素数)
  \end{center}
\end{frame}
%%%%%%%%%%%%%%%

%%%%%%%%%%%%%%%
\begin{frame}{}
  \begin{definition}[最大元与最小元 (Maximum/Minimum; Greatest/Least; Largest/Smallest)]
    令 $\preceq \;\subseteq X \times X$ 是 $X$ 上的偏序。
    设 $a \in X$。\\[5pt]

    \pause
    如果
    \[
      \forall x \in X\; x \preceq a,
    \]
    则称 $a$ 是 $X$ 的\red{最大元}。\\[10pt]

    \pause
    如果
    \[
      \forall x \in X\; a \preceq x,
    \]
    则称 $a$ 是 $X$ 的\red{最小元}。
  \end{definition}

  \pause
  \vspace{0.30cm}
  \begin{center}
    \red{$Q:$ 最大/最小元是否一定存在? 如果存在, 是否唯一?}
  \end{center}
\end{frame}
%%%%%%%%%%%%%%%

%%%%%%%%%%%%%%%
\begin{frame}{}
  \begin{theorem}
    偏序集 $(X, \preceq)$ 如果有最大元或最小元, 则它们是唯一的。
  \end{theorem}

  \pause
  \vspace{0.30cm}
  \begin{center}
    \red{假设存在两个最大元 $x$ 与 $y$。}

    \pause
    \[
      x \preceq y \land y \preceq x \pause \implies x = y
    \]
    \vspace{-0.60cm}
    \begin{center}
      \cyan{(反对称性)}
    \end{center}
  \end{center}
\end{frame}
%%%%%%%%%%%%%%%

%%%%%%%%%%%%%%%
\begin{frame}{}
  \begin{columns}
    \column{0.50\textwidth}
      \fig{width = 0.80\textwidth}{figs/Hasse-Powerset}
    \column{0.50\textwidth}
      \fig{width = 0.80\textwidth}{figs/powerset}
  \end{columns}
\end{frame}
%%%%%%%%%%%%%%%

%%%%%%%%%%%%%%%
\begin{frame}{}
  \begin{definition}[线性拓展 (Linear Extension)]
    设 $(X, \preceq)$ 是偏序集, $(X, \preceq')$ 是全序集。
    如果
    \[
      \forall x, y \in X\; x \preceq y \to x \preceq' y,
    \]
    则称 $\preceq'$ 是 $\preceq$ 的\red{线性拓展}。
  \end{definition}

  \fig{width = 0.40\textwidth}{figs/dag-topo}

  \pause
  \vspace{-0.50cm}
  \[
    \blue{5, 7, 3, 11, 8, 2, 9, 10} \qquad \cyan{3, 5, 7, 8, 11, 2, 9, 10}
  \]
\end{frame}
%%%%%%%%%%%%%%%

%%%%%%%%%%%%%%%
\begin{frame}{}
  \begin{theorem}
    设 $(X, \preceq)$ 是偏序集且 \red{$X$是有限集}, 则 $\preceq$ 的线性拓展必定存在。
  \end{theorem}

  \fig{width = 0.40\textwidth}{figs/dag-topo}

  \pause
  \begin{center}
    \blue{\bf Topological sorting}: iteratively remove minimal elements
  \end{center}

  \pause
  \begin{theorem}
    设 $(X, \preceq)$ 是偏序集且 \red{$X$是有限集}, 则极小元 (极大元) 一定存在。
  \end{theorem}
\end{frame}
%%%%%%%%%%%%%%%

%%%%%%%%%%%%%%%
\begin{frame}{}
  \begin{definition}[上界 (Upper Bound) 与下界 (Lower Bound)]
    设 $(X, \preceq)$ 是\blue{偏序集}。对于 $Y \subseteq X$, 如果
    \[
      \exists \red{x \in X}\; \forall y \in Y\; y \preceq x,
    \]
    则称 $x$ 为 $Y$ 的\red{上界}。

    \pause
    \vspace{0.30cm}
    类似地, 如果
    \[
      \exists \red{x \in X}\; \forall y \in Y\; x \preceq y,
    \]
    则称 $x$ 为 $Y$ 的\red{下界}。
  \end{definition}

  \pause
  \fig{width = 0.55\textwidth}{figs/bound}
\end{frame}
%%%%%%%%%%%%%%%

%%%%%%%%%%%%%%%
\begin{frame}{}
  \begin{definition}[最小上界 (Least Upper Bound; LUB); 上确界 (Supremum; $\sup(Y)$)]
    设 $(X, \preceq)$ 是\blue{偏序集}, $Y \subseteq X$ 是 $X$ 的子集,
    $x \in X$ 是 $Y$ 的\cyan{上界}。 \pause\\[5pt]
    若对于 $Y$ 的所有\cyan{上界} $x'$, 均有 $x \preceq x'$,
    则称 $x$ 是 $Y$ 的\red{最小上界/上确界}。
  \end{definition}

  \pause
  \fig{width = 0.55\textwidth}{figs/bound}

  \pause
  \begin{definition}[最大下界 (Greatest Lower Bound; GLB); 下确界 (Infimum; $\inf(Y)$)]
    设 $(X, \preceq)$ 是\blue{偏序集}, $Y \subseteq X$ 是 $X$ 的子集,
    $x \in X$ 是 $Y$ 的\cyan{下界}。 \\[5pt]
    若对于 $Y$ 的所有\cyan{下界} $x'$, 均有 $x' \preceq x$,
    则称 $x$ 是 $Y$ 的\red{最大下界/下确界}。
  \end{definition}
\end{frame}
%%%%%%%%%%%%%%%

%%%%%%%%%%%%%%%
\begin{frame}{}
  \begin{exampleblock}{}
    \red{如下 $(A, \le)$ 有上确界吗?}
    \[
      A = \set{p \in \Q \mid p > 0 \land p^2 < 2}
    \]
  \end{exampleblock}

  \pause
  \vspace{1em}
  \begin{center}
    \blue{在 $(\R, \le)$ 中}, $A$ \blue{有}上确界: $\sup(A) = \sqrt{2}$

    \pause
    \vspace{1em}
    \red{在 $(\Q, \le)$ 中}, $A$ \red{无}上确界
  \end{center}
\end{frame}
%%%%%%%%%%%%%%%

%%%%%%%%%%%%%%%
\begin{frame}{}
  \begin{theorem}
    集合 $A$ 在 $(\Q, \le)$ 中无上确界。
    \[
      A = \set{p \in \Q \mid p > 0 \land p^2 < 2}
    \]
  \end{theorem}

  \pause
  \begin{proof}
    \[
      B = \set{p \in \Q \mid p > 0 \land p^2 > 2}
    \]

    \pause
    \begin{center}
      只需证明: $A$ 无最大元、$B$ 无最小元
    \end{center}

    \pause
    \[
      \forall p \in A\; \red{\exists q \in A}\; p < q
      \qquad
      \forall p \in B\; \red{\exists q \in B}\; p > q
    \]

    \pause
    \[
      q = p - \frac{p^2 - 2}{p + 2} = \frac{2p + 2}{p + 2}
      \qquad
      q^2 - 2 = \frac{2(p^2 -2)}{(p + 2)^2}
    \]

    \pause
    \[
      p \in A \implies q \in A \land p < q
      \pause \qquad
      p \in B \implies q \in B \land p > q
    \]
  \end{proof}
\end{frame}
%%%%%%%%%%%%%%%

%%%%%%%%%%%%%%%
\begin{frame}{}
  \begin{definition}[实数 ($\R$)]
    如果集合 $A \subseteq \Q$ 满足:
    \begin{enumerate}[(1)]
      \item $A \neq \emptyset$ 且 $A \neq \Q$;
      \item $A$ 是向下封闭的:
        \[
          \forall p \in A, q \in \Q\; (q < p \to q \in A)
        \]
      \item $A$ 没有最大元:
        \[
          \forall p \in A\; \exists q \in A\; p < q
        \]
    \end{enumerate}
    则称 $A$ 是\red{\bf 戴德金分割} (Dedekind cut),
    $A$ 所对应的实数为 $\sup(A)$。
  \end{definition}

  \pause
  \begin{definition}[$\le_{\R}$]
    $x_{1} \le_{\R} x_{2} \iff x_{1} \subseteq x_{2}$
  \end{definition}

  \pause
  \begin{definition}[$+_{\R}$]
    $x +_{\R} y = \set{p +_{\Q} q \mid p \in x, q \in y}$
  \end{definition}
\end{frame}
%%%%%%%%%%%%%%%